% ============================================================================
% MUSTERLÖSUNG: Wiederholung Bewegung (LP-05)
% Physik Klasse 6 - Übungsstunde vor Test
% Stand: 29.01.2026
% ============================================================================

\documentclass[11pt,a4paper]{article}

\usepackage[utf8]{inputenc}
\usepackage[T1]{fontenc}
\usepackage[ngerman]{babel}
\usepackage{geometry}
\usepackage{helvet}
\usepackage{tikz}
\usepackage{enumitem}
\usepackage{fancyhdr}
\usepackage{xcolor}
\usepackage{tcolorbox}
\usepackage{amssymb}
\usepackage{amsmath}
\usepackage{siunitx}

\sisetup{locale=DE, per-mode=symbol}

\renewcommand{\familydefault}{\sfdefault}

\geometry{left=1.5cm, right=1.5cm, top=1.5cm, bottom=1.2cm}

% Fachfarbe Physik
\definecolor{physik}{HTML}{2c5aa0}
\colorlet{fachfarbe}{physik}

% Lösungsfarbe
\definecolor{loesung}{HTML}{c00000}

% Hilfsfarben
\definecolor{gridcolor}{HTML}{c0d0e8}

% ============================================================================
% KOPFZEILE
% ============================================================================
\pagestyle{fancy}
\fancyhf{}
\fancyhead[L]{\small\textbf{Physik Klasse 6} | Wiederholung Bewegung}
\fancyhead[R]{\small\textcolor{loesung}{\textbf{Musterlösung}}}
\renewcommand{\headrulewidth}{0.4pt}

% ============================================================================
% MERKKASTEN
% ============================================================================
\tcbuselibrary{skins}

\newtcolorbox{merkkasten}[1][]{
    colback=white,
    colframe=fachfarbe,
    coltitle=black,
    colbacktitle=white,
    fonttitle=\bfseries\small,
    boxrule=1pt,
    arc=2pt,
    left=8pt, right=8pt, top=6pt, bottom=6pt,
    #1
}

% ============================================================================
% AUFGABENTITEL
% ============================================================================
\newcommand{\aufgabe}[1]{\noindent\textbf{\textcolor{fachfarbe}{#1}}}
\newcommand{\lsg}[1]{\textcolor{loesung}{\textbf{#1}}}

\begin{document}

% ============================================================================
% TITEL
% ============================================================================
\noindent
{\Large\bfseries Wiederholung: Bewegung} \hfill \textcolor{loesung}{\Large\textbf{MUSTERLÖSUNG}}
\vspace{0.3em}
\hrule
\vspace{0.6em}

% ============================================================================
% BLOCK A: BEWEGUNGSARTEN
% ============================================================================
\aufgabe{Aufgabe A: Kreuze die richtige Bahnform und Bewegungsart an.}

\vspace{0.3em}
{\small
\begin{tabular}{|p{5.5cm}|c|c|c|c|}
\hline
\textbf{Beispiel} & \multicolumn{2}{c|}{\textbf{Bahnform}} & \multicolumn{2}{c|}{\textbf{Bewegungsart}} \\
\cline{2-5}
 & geradl. & krumml. & gleichf. & ungleichf. \\
\hline
a) ICE fährt mit 200\,km/h geradeaus & \lsg{$\boxtimes$} & $\square$ & \lsg{$\boxtimes$} & $\square$ \\
\hline
b) Auto bremst an der Ampel & \lsg{$\boxtimes$} & $\square$ & $\square$ & \lsg{$\boxtimes$} \\
\hline
c) Karussell dreht gleichmäßig & $\square$ & \lsg{$\boxtimes$} & \lsg{$\boxtimes$} & $\square$ \\
\hline
d) Fahrrad fährt eine Kurve & $\square$ & \lsg{$\boxtimes$} & $\square$ & \lsg{$\boxtimes$} \\
\hline
e) Rolltreppe fährt gleichmäßig hoch & \lsg{$\boxtimes$} & $\square$ & \lsg{$\boxtimes$} & $\square$ \\
\hline
f) Achterbahn fährt durch Looping & $\square$ & \lsg{$\boxtimes$} & $\square$ & \lsg{$\boxtimes$} \\
\hline
\end{tabular}
}

\vspace{0.6em}

% ============================================================================
% BLOCK B: GESCHWINDIGKEIT
% ============================================================================
\begin{merkkasten}[title=Formel: Geschwindigkeit]
\begin{center}
{\large $v = \dfrac{s}{t}$} \qquad umgestellt: \quad $t = \dfrac{s}{v}$ \qquad $s = v \cdot t$
\end{center}
\end{merkkasten}

\vspace{0.4em}

\aufgabe{Aufgabe B: Berechne die fehlende Größe.}

\vspace{0.3em}
{\small
\begin{tabular}{|c|p{6cm}|p{8cm}|}
\hline
& \textbf{Angaben} & \textbf{Rechnung} \\
\hline
B1 & $s = 200$\,m, $t = 40$\,s & \lsg{$v = \dfrac{s}{t} = \dfrac{200\,\text{m}}{40\,\text{s}} = 5\,\si{\metre\per\second}$} \\[0.6cm]
\hline
B2 & $v = 4$\,\si{\metre\per\second}, $s = 200$\,m & \lsg{$t = \dfrac{s}{v} = \dfrac{200\,\text{m}}{4\,\si{\metre\per\second}} = 50\,\text{s}$} \\[0.6cm]
\hline
B3 & $v = 6$\,\si{\metre\per\second}, $t = 30$\,s & \lsg{$s = v \cdot t = 6\,\si{\metre\per\second} \cdot 30\,\text{s} = 180\,\text{m}$} \\[0.6cm]
\hline
B4 & $s = 150$\,m, $t = 30$\,s & \lsg{$v = \dfrac{s}{t} = \dfrac{150\,\text{m}}{30\,\text{s}} = 5\,\si{\metre\per\second}$} \\[0.6cm]
\hline
B5 & $v = 8$\,\si{\metre\per\second}, $s = 160$\,m & \lsg{$t = \dfrac{s}{v} = \dfrac{160\,\text{m}}{8\,\si{\metre\per\second}} = 20\,\text{s}$} \\[0.6cm]
\hline
B6 & $v = 5$\,\si{\metre\per\second}, $t = 40$\,s & \lsg{$s = v \cdot t = 5\,\si{\metre\per\second} \cdot 40\,\text{s} = 200\,\text{m}$} \\[0.6cm]
\hline
\end{tabular}
}

\vspace{0.6em}

% ============================================================================
% BLOCK C: s-t-DIAGRAMME
% ============================================================================
\aufgabe{Aufgabe C1: Lies die Werte aus dem Diagramm ab.}

\vspace{0.3em}
\begin{minipage}{0.55\textwidth}
\begin{tikzpicture}[scale=0.65]
    \draw[gridcolor, very thin] (0,0) grid (8,5);
    \draw[->, thick] (0,0) -- (8.5,0) node[right] {$t$ in s};
    \draw[->, thick] (0,0) -- (0,5.5) node[above] {$s$ in m};
    \foreach \x in {0,2,4,6,8} { \node[below] at (\x,0) {\small \x}; }
    \foreach \y/\label in {0/0,1/20,2/40,3/60,4/80} { \node[left] at (0,\y) {\small \label}; }
    \draw[fachfarbe, very thick] (0,0) -- (8,4);
    % Markierungen bei t=4 und t=6
    \draw[fachfarbe, thick] (3.85,1.85) -- (4.15,2.15);
    \draw[fachfarbe, thick] (3.85,2.15) -- (4.15,1.85);
    \draw[fachfarbe, thick] (5.85,2.85) -- (6.15,3.15);
    \draw[fachfarbe, thick] (5.85,3.15) -- (6.15,2.85);
    % Hilfslinien für Ablesen
    \draw[loesung, dashed] (4,0) -- (4,2) -- (0,2);
    \draw[loesung, dashed] (6,0) -- (6,3) -- (0,3);
\end{tikzpicture}
\end{minipage}
\hfill
\begin{minipage}{0.42\textwidth}
a) Weg nach $t = 4$\,s: \quad $s$ = \lsg{40\,m}

\vspace{0.6em}
b) Zeit bei $s = 60$\,m: \quad $t$ = \lsg{6\,s}

\vspace{0.6em}
c) Weg nach $t = 2$\,s: \quad $s$ = \lsg{20\,m}

\vspace{0.6em}
d) Zeit bei $s = 80$\,m: \quad $t$ = \lsg{8\,s}
\end{minipage}

\newpage

% ============================================================================
% SEITE 2
% ============================================================================

\aufgabe{Aufgabe C2: Vergleiche die Läufer und begründe.}

\vspace{0.5em}
\begin{minipage}{0.5\textwidth}
\begin{tikzpicture}[scale=0.8]
    \draw[gridcolor, very thin] (0,0) grid (8,5);
    \draw[->, thick] (0,0) -- (8.5,0) node[right] {$t$ in s};
    \draw[->, thick] (0,0) -- (0,5.5) node[above] {$s$ in m};
    \foreach \x in {0,2,4,6,8} { \node[below] at (\x,0) {\small \x}; }
    \foreach \y/\label in {0/0,1/20,2/40,3/60,4/80} { \node[left] at (0,\y) {\small \label}; }
    \draw[red, very thick] (0,0) -- (4,4);
    \node[red] at (4.5,4) {\small A};
    \draw[blue, very thick] (0,0) -- (8,4);
    \node[blue] at (8.3,4) {\small B};
\end{tikzpicture}
\end{minipage}
\hfill
\begin{minipage}{0.47\textwidth}
a) Kreuze an: Wer ist schneller?

\vspace{0.3em}
\quad \lsg{$\boxtimes$} Läufer A \quad $\square$ Läufer B

\vspace{0.8em}
b) Begründe deine Antwort:

\vspace{0.2em}
\lsg{Läufer A legt die gleiche Strecke (80\,m) in kürzerer Zeit (4\,s statt 8\,s) zurück. Die Linie von A ist steiler.}

\vspace{0.8em}
c) Berechne $v$ für Läufer A:

\vspace{0.2em}
\lsg{$v_A = \dfrac{s}{t} = \dfrac{80\,\text{m}}{4\,\text{s}} = 20\,\si{\metre\per\second}$}

\vspace{0.8em}
d) Berechne $v$ für Läufer B:

\vspace{0.2em}
\lsg{$v_B = \dfrac{s}{t} = \dfrac{80\,\text{m}}{8\,\text{s}} = 10\,\si{\metre\per\second}$}
\end{minipage}

\vspace{1em}

% ============================================================================
% BLOCK D: VERKEHRSSICHERHEIT
% ============================================================================
\begin{merkkasten}[title=Verkehrssicherheit]
\textbf{Anhalteweg} = Reaktionsweg + Bremsweg \quad | \quad Bei 1\,s Reaktionszeit: $s_R$ = $v$
\end{merkkasten}

\vspace{0.4em}

\aufgabe{Aufgabe D: Berechne Reaktionsweg und Anhalteweg.}

\vspace{0.2em}
{\small
\begin{tabular}{|c|c|c|c|c|}
\hline
& \textbf{Geschwindigkeit} & \textbf{Reaktionsweg} & \textbf{Bremsweg} & \textbf{Anhalteweg} \\
\hline
D1 & 5\,\si{\metre\per\second} & \lsg{5\,m} & 2,5\,m & \lsg{7,5\,m} \\
\hline
D2 & 10\,\si{\metre\per\second} & \lsg{10\,m} & 10\,m & \lsg{20\,m} \\
\hline
D3 & 8\,\si{\metre\per\second} & \lsg{8\,m} & 6,4\,m & \lsg{14,4\,m} \\
\hline
D4 & 15\,\si{\metre\per\second} & \lsg{15\,m} & 22,5\,m & \lsg{37,5\,m} \\
\hline
D5 & 12\,\si{\metre\per\second} & \lsg{12\,m} & 14,4\,m & \lsg{26,4\,m} \\
\hline
\end{tabular}
}

\vspace{0.6em}

\aufgabe{Aufgabe E: Erkläre, warum der Sicherheitsabstand bei höherer Geschwindigkeit größer sein muss.}

\vspace{0.3em}
\lsg{Bei höherer Geschwindigkeit ist sowohl der Reaktionsweg als auch der Bremsweg größer. Der Anhalteweg wächst also mit der Geschwindigkeit. Um bei einer Gefahrenbremsung nicht aufzufahren, muss der Sicherheitsabstand mindestens so groß wie der Anhalteweg sein.}

\end{document}
